% Chinese review manuscript. Compile with XeLaTeX -> BibTeX -> XeLaTeX x2.
% CSCWD requires an English, IEEE proceedings paper of at most six pages.
% This manuscript reports implementation evidence; participant results do not exist.
\RequirePackage[T1]{fontenc}
\documentclass[conference]{IEEEtran}
\usepackage[UTF8,scheme=plain,fontset=none]{ctex}
\setmainfont{TeX Gyre Termes}
\setsansfont{TeX Gyre Heros}
\setmonofont{Latin Modern Mono}
\setCJKmainfont{FandolSong-Regular}[Script=Default,BoldFont=FandolSong-Bold,ItalicFont=FandolKai-Regular,BoldItalicFont=FandolKai-Regular,SmallCapsFont=FandolSong-Regular]
\setCJKsansfont{FandolHei-Regular}[Script=Default,BoldFont=FandolHei-Bold]
\setCJKmonofont{FandolFang-Regular}[Script=Default]
\usepackage{amsmath,amssymb}
\usepackage{graphicx}
\usepackage{booktabs,tabularx,array}
\usepackage{balance}
\usepackage{cite}
\usepackage[hidelinks]{hyperref}
\usepackage{url}
\newcolumntype{Y}{>{\raggedright\arraybackslash}X}
\graphicspath{{figures/}}
\newcommand{\sys}{CoVisStudio}
\newcommand{\Bzero}{B0}
\newcommand{\Btwo}{B2}
\renewcommand{\abstractname}{摘要}
\renewcommand{\IEEEkeywordsname}{关键词}
\renewcommand{\refname}{参考文献}
\renewcommand{\figurename}{图}
\renewcommand{\tablename}{表}
\hyphenation{CoVisStudio}

\title{CoVisStudio：面向全双工协同可视分析的\\证据与修订范围分离}
\author{\IEEEauthorblockN{匿名作者：中文审阅稿}
\IEEEauthorblockA{方法、系统实现与初步技术验证\\正式用户对照实验尚未开展}}

\begin{document}
\maketitle

\begin{abstract}
用户与人工智能共同修改分析看板时，可能在助手讲解期间改变需求。例如，用户要求复用刚才讨论的三个品类，但继续计算当前地区的指标。此时，引用历史对象并不意味着恢复历史筛选条件。全双工交互使用户可以在播放和执行期间发话，也使生成内容、已播放内容及当前画面可能不一致。本文提出并实现 \sys，将一次修订明确表示为对象依据、修改目标与内容、保留与解除条件，并通过可纠正的修订卡支持双方核对。系统共用历史集合工具和提交规则，分别记录图表提交、渲染、可见性及语音播放。初步技术验证表明，本地脚本驱动的浏览器流程可完成历史集合复用及播放中日期修订，24项机制检查通过。早期两个真实案例均未完成全部目标，随后八次诊断四次最终状态正确；本轮另两次复测一次成功、一次仍因参数类型失败。新记录验证晚到转写关联，但未触发旧的迟到响应竞态；提交后仍上传旧图像的缺口随后仅经本地修复。本文据此给出仅含两个全双工条件的对照协议，用于检验修订正确性、用户理解与修复成本。现有证据支持系统机制的可实现性，尚不支持用户收益或全双工优越性的结论。
\end{abstract}
\begin{IEEEkeywords}
人机协作；全双工语音；可视分析；意图修订；共享证据
\end{IEEEkeywords}

\section{引言}
分析看板的制作通常不是一次性地执行完整需求。用户先观察图表，再选择比较对象、调整指标或缩小范围。人工智能参与这一过程时，双方需要持续确认正在讨论什么、哪些要求已经改变，以及哪些条件仍然有效。这是一种围绕共同分析产物的协作，而不只是把自然语言翻译成绘图命令。

考虑一个包含订单量、低评分率和配送时长的三图看板。用户先在 RJ 州选出三个值得关注的品类，随后转到 MG 州。助手讲解期间，用户说：“就用你刚才说的那三个品类，修改右边两张服务图，地区和日期保持现在的。”正确操作需要复用历史品类成员、修改指定视图，并保留当前 MG 范围。如果系统恢复整份历史查询，或重新选择 MG 的前三品类，图表可能仍然数值正确，却回答了另一个分析问题。

全双工语音允许用户在助手播放时继续发话。它为这种及时修订提供了入口，但并不自动保证双方理解一致。模型已生成的句子可能尚未播放；图表已提交却未完成渲染；用户引用的讲解可能对应历史对象而非最新集合。因此，单独保存完整聊天记录或最后一张图，不能完整说明一次修订的依据与过程。

已有研究已经支持自然语言歧义处理、语音与画布协作，以及执行中的用户介入\cite{datatone,cleo,voca}。本文不将“可以打断”或“语音加图片”作为新贡献。我们聚焦以下问题：\emph{在允许用户边听边修订的协作中，将引用证据与当前修改范围显式分开，能否帮助双方改对并理解改动，其交互代价是什么？}两种实验条件均使用全双工环境；本文不比较全双工与顺序交互。

本文工作包括两部分。第一，提出并实现面向分析修订的三部分表示，将历史对象身份与当前筛选、指标和目标视图区分，使用可纠正的修订卡使其可见。第二，建立连接输入、实际执行与结果交付的验证流程，报告已完成的技术检查及真实失败，并给出不削弱基线的用户对照协议。用户效果是该协议要检验的假设，不是本文已有的实证结论。

\section{相关工作}
\subsection{自然语言分析与歧义处理}
DataTone 将自然语言歧义显式呈现，允许用户纠正解释，并以部分约束维持跨查询纠正\cite{datatone}。这与本文“让解释可检查”的思想直接相关。本文的区别在于具体修订对象：用户可以从历史讲解引用成员，同时要求保留当前计算条件。我们将对象来源与应用范围分开，并记录与修订有关的实际播放和可见版本。该区别是方法关注点的变化；不能仅凭使用新语音接口就宣称超越 DataTone。

Speech-to-Visualization 研究从语音问题与数据库信息生成可视化查询\cite{speechvis}。ChartEditBench 评估基于现有图表与代码的多轮编辑\cite{charteditbench}。这些工作构成语音生成与改图任务的直接基础。本文关注的不是新增图形类型，而是复用历史证据时，比较对象和当前分析条件是否仍符合用户要求。

\subsection{语音与共享画布协作}
CommSketch 结合草图和口述帮助传达设计意图\cite{commsketch}。CLEO 在共享设计画布中跟踪用户动作并选择性调整计划，研究执行中的协作\cite{cleo}。VoCa 协调语音、画布创建、注释与注意引导，并以探索性部署研究使用体验\cite{voca}。因此，边说边做、共享画面和中途修订已有明确先例。

本文进一步选择具有可核验数据口径的分析看板，研究“历史对象仍可引用，但旧范围不应自动恢复”的修订。与开放画布相比，该场景能够分别检查成员集合、筛选范围、指标口径与未修改视图。它提供更明确的结果核验方式，但场景变化本身不足以证明研究贡献。需要通过公平对照验证显式支持是否改善理解与修订，而不仅是让界面看起来更完整。

\subsection{全双工工具执行与共同理解}
Full-Duplex-Bench-v3 已研究真实非流畅语音下的工具使用\cite{fdbv3}；DuplexAgent-RSI 已处理实时语音与后台执行的协调，包括失效结果的交付管理\cite{duplexagent}。本文沿用中断与提交管理作为工程基础，不将其声明为新算法。

共同理解依赖双方持续提供并核对理解证据，而非仅发出话语\cite{grounding}。据此，本文将修订卡设计为可纠正的共享表示。模型内部选中了正确集合，与用户知道选中了哪个集合，是两个不同结果，必须分别测量。

\section{问题定义与修订方法}
\subsection{引用证据与当前状态}
以一次用户发话及随后澄清、执行和修复构成一个修订回合。回合不等于单个模型响应：一个请求可能产生工具调用、后续语音和多个交付事件。

对当前视图 $i$，记其查询配置为 $C_i$，包含地区、日期、订单状态、指标及成员选择方式。历史证据 $e$ 包含来源版本和当时的成员集合 $M(e)$。一次修订表示为
\begin{equation}
R=(O,T,\Delta,P,L),
\end{equation}
其中 $O$ 是引用对象的来源，$T$ 是当前目标视图，$\Delta$ 是明确要求的修改，$P$ 是保留条件，$L$ 是明确解除的条件。界面将其组织为三部分：\emph{对象依据；目标与修改；保留与解除}。这些字段来自用户表达及当前证据，不包含评分答案。

例如，上述请求选择历史集合 $M(e)$，目标为图2和图3，保留 MG、当前日期、已送达订单和各图指标。图1保持不变。若用户随后要求“不要固定这三个，重新找 MG 的前三”，则解除成员固定约定并恢复明确的排名规则。固定成员与保留动态排名规则必须分开表示。

\subsection{修订的执行规则}
系统先列出真实存在的历史图表和已保存集合，模型从中引用候选标识，而非自由编造成员来源。用户也可在界面直接选定证据。两者使用同一成员解析器。引用历史集合时，只复制对象身份，不复制旧地区、日期和指标。

对目标视图，执行器在当前配置上施加明确修改；未指定的字段保留。选择历史对象时，可写为
\begin{equation}
C_i'=\operatorname{patch}(C_i,\Delta_i,M(O)),\quad i\in T.
\end{equation}
对 $i\notin T$，有 $C_i'=C_i$。若选择动态重排，则必须给出排名指标和数量，并解除旧固定成员。以上规则检查操作是否符合已表达的解释，不能替代模型理解含糊的话语。

明确请求通过单次 \texttt{revise\_analysis} 调用完成私有计算与提交。工具只接收平坦参数，如目标视图、成员动作、候选标识和需要改变的查询字段。B0与B2共用成员绑定、编译器及真实回执；B2另显示摘要及可展开卡片，B0保留直接呈现。卡片不要求每次点击确认：只有尚未解决的选择才暂停提交并澄清。明确请求通常先提交再显示卡片，支持检查与后续纠正，不是一律提交前验证。

这使卡片既不成为所有操作的额外门槛，也不被当作“用户已经同意”的证据。卡片生成、进入视口、展开查看和成功确认分别记录。用户看到了卡片，不等于读懂了它。

\begin{figure*}[t]
\centering
\includegraphics[width=\textwidth]{fig01_revision_timeline.pdf}
\caption{全双工修订的示意场景。历史集合可作为对象依据，当前范围仍为 MG；模型生成的后续讲解不等于用户已经听到。示意顺序不代表真实实验时长，集合符号不代表测量数据。}
\label{fig:timeline}
\end{figure*}

\subsection{全双工在方法中的作用}
全双工让用户在播放期间表达新的要求，因而修订可能发生在解释尚未交付完、工具尚未完成或图表刚刚更新的时候。图\ref{fig:timeline}展示的困难不是单纯“旧结果覆盖新结果”，还包括解释正确执行却引用了错误对象的语义错误。

收到语音开始事件后，系统立即撤销旧工作的后续发布资格并停止旧播放，再请求远端取消。旧工作此前已经合法提交的结果仍保留，不被追溯撤销。本文的执行保证始于本地处理事件，不能追回已送达的音频，也不能等同于零延迟识别用户的物理发声。

“用 RJ 保存的集合”可由显式名称解析，不一定需要播放位置；“就用你刚才说的那组”才可能依赖实际讲解片段。是否真的存在这种依赖，需要人工结合话语与交付证据判断。不同候选产生不同图表，只能说明选择有后果，不能单独证明全双工信息必不可少。

\section{系统实现与证据记录}
\subsection{界面、模型与工具}
\sys 使用 Vue 和 ECharts 实现英文三图界面，后端使用 FastAPI 管理会话、工具和公开图表状态。语音模型为 Qwen3.8-Omni-Flash-Realtime\cite{qwenrealtime}。界面围绕图表、对话、保存集合和修订卡组织，支持选中当前视图及检查历史版本。语音连接只在显式开始后建立。

图表数据来自 Olist\cite{olist}。订单量、低评分率和配送时长由真实数据查询得到。低评分率以具有有效评分的订单为分母，配送时长使用有效配送记录；实现保留各指标口径。跨品类订单可能重复参与不同品类的比较，因此不把品类计数简单相加当作唯一订单总数。

工具在私有候选中计算，提交时同时检查请求是否仍有效、当前图表版本是否匹配。两种条件共用候选解析、已保存集合工具、数据查询与提交规则。语音回复中的成功表述不能触发提交。系统根据当前请求实际产生的事件显示“已提交”“已显示”或“未修改图表”，上一轮的成功回执不会证明本轮完成。

\begin{figure*}[t]
\centering
\includegraphics[width=\textwidth]{fig02_revision_architecture.pdf}
\caption{系统中的修订与验证链。B0 与 B2 共用证据读取、工具执行和提交规则；B2 向用户显式呈现可纠正的修订结构。图表提交、实际渲染、可见性与音频播放分别记录。应用证据包不等于冻结模型原生会话历史。}
\label{fig:architecture}
\end{figure*}

\subsection{证据截止与实际交付}
系统用请求、响应、工具和图表版本标识关联事件。浏览器与服务器通过时间交换估计时钟偏移及不确定度，同时保留原始时钟。每条应用证据具有采集时间 $t_c$、服务器可用时间 $t_a$、来源和内容摘要。准备一次解释时，以截止时间 $d$ 选择证据：
\begin{equation}
t_c+u\leq d\quad\land\quad t_a\leq d,
\end{equation}
其中 $u$ 是采集时间的不确定度。截图虽然早已采集，但尚未到达服务器时，不计为该时刻可用的应用输入。无法确定时间关系的材料保留未知状态。

播放记录由 AudioWorklet 的输出游标生成，区分收到的采样、已输出采样和中断后未播放采样。以24 kHz音频计，约每1,200个源采样产生一次进度回执。完整接收的 WAV 不能作为完整播放的证据；软件输出采样也不能证明用户听懂了哪些词。图表则分别记录提交、绘制完成、进入视口和离开视口，避免把后台提交等同于用户接触到结果。

应用证据包保存版本和摘要，并归档实际出站模型输入。当前实现没有隔离 Qwen 原生会话中的既有内容，因此仅能声明应用层读取和上传受截止规则约束，不能声明模型的全部知识已被冻结。语音与语义片段的对齐仍需单独标注。图\ref{fig:architecture}明确了这些实现边界。

\section{初步技术验证}
本节报告截至2026年11月13日已保存的技术记录。局部脚本驱动、真实模型调用和用户实验是不同证据来源；本文不将其合并计算成功率。两个早期案例、随后本地流程、八次诊断与本轮两次复测分开报告，不合并为十次同版本效果实验。

\subsection{真实模型案例暴露的问题}
两个案例均使用真实 Olist 查询、浏览器和 Qwen Realtime 连接。用户输入由现代语音合成模型生成；三段源语音分别为6.64、8.56和6.32秒。它们是编写的测试请求，不是参与者自然使用数据，也未经过独立核听评分。

案例一复用历史品类，但将嵌套证据对象写成 JSON 字符串，旧流程错误后仍续答，发生24次失败工具尝试而无提交。案例二先改为上半年，播放中要求第一季度，虽停播并正确转录，却只口头承诺，未再改图。

两个案例的六张最终图均匹配其实际查询的独立 SQL，却未满足请求目标。这说明数值正确、回复合理和任务正确必须分开检查，未重新连接也不等于未重复调用。

\subsection{修改后的本地流程}
针对上述问题，我们改为平坦参数和单次修订调用，停止错误回执后的自动续答，并将实际执行回执关联到当前请求。验证使用本地脚本模型与测试音频，通过真实浏览器、音频工作线程、Olist 工具和图表渲染运行。

流程先复用历史集合并改为上半年，在实际播放期间输入第一季度修订，最后注入一次没有工具调用的口头承诺。公开版本依次为 $6\rightarrow7\rightarrow8\rightarrow8$。图2和图3在第二次提交中只改变日期，保留 MG、历史成员、订单状态和各自指标；图1不变。最后一次承诺不产生新版本，界面显示“未修改图表”，实际界面见图\ref{fig:workspace}。

\begin{figure*}[t]
\centering
\includegraphics[width=\textwidth]{fig03_actual_workspace.pdf}
\caption{本地脚本验证的真实浏览器画面，分区裁剪自同一截图。图1仍是全年订单量，图2和图3为第一季度；历史 RJ 成员依据独立显示。随后请求第二季度却没有工具调用，因此界面保留第一季度并提示“未修改图表”。这不是用户实验或真实模型成功截图。}
\label{fig:workspace}
\end{figure*}

\begin{table}[t]
\caption{已完成验证及其证据边界}
\label{tab:evidence}
\centering\small
\begin{tabularx}{\columnwidth}{@{}p{.21\columnwidth}YY@{}}
\toprule
记录 & 实际观察 & 不能据此推断 \\
\midrule
真实案例一 & 参数格式错误；无提交 & 所有引用任务都会失败 \\
真实案例二 & 插话停播；未执行日期修订 & 失败由全双工本身导致 \\
本地集成 & 24项检查通过；两次提交、一次不提交 & 真实模型已稳定改对 \\
用户实验 & 尚未开展 & B2 提高正确率或减轻负担 \\
\bottomrule
\end{tabularx}
\end{table}

此前该流程的24项机制检查全部通过，涵盖单次修改、重叠停播、独立 SQL、未修改视图保留、图片输入、版本匹配的渲染及关闭。24是同一流程的断言数量，不是24个独立任务。两种路径的共用成员解析与执行规则也通过局部检查。这支持执行链可运行，不证明真实模型会选对工具或理解指代。

\subsection{修改后的真实模型诊断}
我们冻结四类编写请求，在运行时 \texttt{direct} 与 \texttt{bound\_revision} 路径各执行一次，共八次真实 Qwen Realtime 诊断。每次使用真实浏览器和 Olist，输入源语音为4.24--8.56秒，连接均少于60秒，无自动重试。修订输入均与实际助手播放重叠。历史引用任务预先选择了证据，因此检验的是成员迁移与当前范围保持，不能证明模型依据实际听到的片段消解指代。评分依据测试前确定的请求，检查最终查询配置、未修改视图及独立 SQL；这些答案不进入模型输入。

\begin{table}[t]
\caption{八次编写请求的真实模型诊断：最终状态}
\label{tab:real-matrix}
\centering\small
\begin{tabularx}{\columnwidth}{@{}p{.27\columnwidth}YY@{}}
\toprule
任务 & direct & bound\_revision \\
\midrule
历史集合及季度修订 & 正确；关闭超时标记 & 仍为上半年；未调用工具 \\
解除固定并本地重排 & 参数类型拒绝 & 参数类型拒绝 \\
历史集合迁移 & 正确 & 取消工具回合后接口错误 \\
保留集合、替换指标 & 正确 & 正确 \\
\bottomrule
\end{tabularx}
\end{table}

八次中四次最终状态符合预设请求，并有实际提交、匹配渲染和末轮音频自然播放结束记录。原始驱动仅记三次通过：历史集合及季度修订的 direct 尝试已在54.634秒完成末轮播放，但关闭过程中55秒看门狗触发，55.752秒关闭，故原始记录保留失败。这不是目标未完成，也不应静默改写为驱动通过。

其余失败有不同来源。季度修订的 bound 路径正确转录请求，却未调用工具，口头宣称已更新并报告旧上半年数值，长语音随后被预算截断。重排任务两路径均将整数参数 \texttt{top\_n} 写为字符串 \texttt{'3'}，一次工具拒绝后终止。历史迁移的 bound 路径在第二个 VAD 片段撤销旧回合后收到“无法继续已取消工具回合”的真实接口错误；不能归为模型单纯不执行，原始 \texttt{finish\_reason=completed} 也不能证明任务完成。

若干尝试还出现首段转录因 epoch 前进而未进入规范化应用记录，原生会话仍保留更多上下文。这暴露了应用记录与模型历史的边界。上述运行路径尚未达到正式 B0/B2 输入等价性冻结要求；四次正确仅描述本矩阵，不估计总体成功率或方法收益。收到音频、软件自然播放结束与人实际听懂仍不同，独立核听和用户理解均待验证。

\subsection{本轮两次独立复测}
在后续修复片段关联与迟到响应流程后，另行授权复测旧矩阵的02-direct与03-bound。它们复用完整源音频，不新增合成或自动重试，连接分别19.890与33.750秒；解释源4.24秒，两段修订源8.08与6.64秒，均未截断且不足10秒。

02仍把\texttt{top\_n}写为字符串\texttt{"3"}，工具拒绝后停止，三图保持r6旧固定集合。源文本中的“Release fixed categories”被识别为“These fixed categories”；前半句本地前三要求仍被保留，模型也选择了重排。03调用\texttt{revise\_analysis}复用选中历史成员，一次提交r7，只改图2和图3，保留RJ、第三季度和各指标，图1完整不变。六张实际最终图均通过原始数据独立SQL，但02的旧查询不满足请求，03满足。

两次均出现第二次语音开始后首片转写才到达，分别晚97.912与55.335毫秒；两片均合并到最新请求，支持该关联修复在真实时序中生效。03最终音频355,200个24 kHz采样全部播放，即14.8秒，未中断；02仅播放63,360/261,120个采样后停止。软件交付不赋予人工核听或理解通过。

新03没有旧的取消工具回合错误，但也未触发迟到创建的旧工具回合，不能据此验收该竞态。其r7提交后续答仍上传r6图像，仅有r7本地采集与真实工具回执，暴露了最新图像交付缺口；随后补交当前图像的修复只经本地检查，未再付费复测。两条新观察分别为一成一败，保留原八次失败，不作为组间效应证据。

\section{双条件用户实验协议}
本节是用于检验方法价值的\emph{计划协议}，不是已完成的实验结果。研究只保留 B0 与 B2，不增加模态、反馈节奏或人为延迟等实验因素。

\subsection{条件与业务任务}
B0 是具有完整历史读取、保存集合及精确复用能力的强基线；B2 在相同能力上，将对象来源、目标与修改、保留与解除条件以修订卡显式呈现。两组固定模型版本、图像更新、工具、语音中断行为、澄清规则和预算上限。明确请求无需额外确认，普通执行回执两组均可见。比较估计整个显式支持方案的效果，不能把收益单独归因于卡片的某个组件。

当前两路径已共享指令、成员编译器和执行能力，本地应用输入、读取及回执等价性审计通过，原生会话策略与实际呈现仍需正式冻结。比较的是可见、可纠正的支持包，不是绑定算法；不能让一组额外获得成员答案。

任务沿用四类业务材料：地区服务评估、本地优先级重选、季节比较和指标替换。参与者需要说明哪些品类值得关注及依据，但不给出唯一正确的业务建议。任务说明只限定业务目的及必须保持的分析条件；参与者可以保留旧集合、重新选择或解除固定。当前含预设成员答案的诊断脚本不直接充当自然使用任务。

每位参与者完成四次任务，每种条件两次，每个业务家族一次；不同平行版本、条件顺序和任务顺序交叉分配。现有30人招募目标属于计划，并非已招募样本或功效保证。30个分配槽不是完整平衡区组，需保留家族、版本与顺序信息。每次任务上限480秒，其中统一安排一次不超过60秒的语音区间；准备和事后理解题不占该语音区间。不得将其描述为八分钟持续语音。参与者无需按脚本在指定时刻打断。

\subsection{独立意图标注与指标}
两名标注者先依据用户原话、明确选择和当时可用证据，独立确定引用成员、目标视图、改变、保留与解除条件。意图标注阶段隐藏模型提案与最终输出，锁定意图规格后再核对结果；必要的既有讲解与用户当时画面仍可查看。报告标注一致率和分歧，仍不能确定的标为未知。独立 SQL 在意图确定后计算参考结果，评分材料不进入模型输入。

主要指标为最终修订正确性和用户理解。正确性要求最终状态符合最后明确意图、保留有效条件，并通过数据核验；未完成但意图明确的任务可判不正确。理解题分别询问对象来源、影响视图和当前计算范围，每项按预先编码规则记0或1，取三项均值。缺失或无法判定的回答单列未知，不根据界面是否打开自动赋分。

次要指标包括经人工确认的误解修复次数、修复耗时和总任务时间。正常探索、主动改变目标和未完成的中间状态不自动算错误。输出解释是否与实际提交一致单独编码，不能由“音频播放完毕”代替语义判断。卡片查看和确认仅作为过程描述。

\subsection{统计与判定范围}
所有分配任务进入主分析，包括未发生插话、失败、超时和意图未知的任务。各分配任务首次正式开始的尝试为主记录；重复尝试另行保留，不按最新或最佳结果替换。当前旧导出程序的“最新尝试”汇总不能直接用于这一分析。同一人的四次任务不视为独立参与者。

首先报告各条件的已知成功、失败、未知数及效应识别范围，以参与者整簇重采样10,000次给出95\%置信区间。无未知时，每人先计算两条件各两任务的平均得分，再对差值进行双侧配对 $t$ 检验；两个主要指标以 Holm 法控制整体 $\alpha=0.05$。存在未知时，完整病例比较只作敏感性描述，不承担确认性结论。家族、版本与顺序的调整采用参与者固定效应回归，重采样仍以参与者为单位，不临时挑选显著模型。

对未知结果，不默认为失败，也不静默删除。设条件 $c$ 的分配任务数为 $n_c$，已知成功数为 $s_c$，未知数为 $u_c$，则成功比例的识别范围为
\begin{equation}
[L_c,U_c]=[s_c/n_c,(s_c+u_c)/n_c].
\end{equation}
条件差的范围为 $[L_2-U_0,U_2-L_0]$，并按参与者聚类报告不确定性。理解题可同样按未知项的0--1范围进行敏感性分析。完整结果不足时，不以排除未知后的显著性作为确认结论。

时间分析保留失败与超时，分别报告正确完成比例和所有分配任务的观察时间。未达到正确状态的任务在终止时记为删失；观察时间不解释为正确完成速度。计划招募规模需要在不查看条件收益的可行性检查后，用主要指标的变异与精度目标确认。语音预算、标注规则和材料适用性均在采集前冻结。

\section{讨论与限制}
\subsection{能够支持的贡献}
本文将“从何处引用对象”与“对当前分析做什么”分开，并把解释结果与实际提交对应起来。对协同设计而言，用户不是批准一段模糊的“已修改”回复，而是能够核对共同产物的对象、范围和变化。这一方法应由修订正确性、理解与代价共同评价。

现阶段可以确认的是方法和记录链已实现，局部诊断能复现并区分参数失败、未执行承诺和正确提交。不能确认的是显式支持是否比强基线更有效。若未来仅理解得分提高，结论应限于改动理解；若正确性提高但时间增加，应报告权衡；若区间很宽，则不能以“没有显著差异”主张等效。

\subsection{全双工与新颖性的边界}
两组都支持实时插话，因此实验不检验全双工是否优于文本或轮次语音。全双工的作用是研究环境：它允许用户在部分讲解已交付、部分工作未结束时修订，为对象引用和当前状态的分离提供具体场景。若自然任务中极少出现依赖实际播放片段的引用，则论文只能讨论一般多轮分析修订的支持，不能将结果扩大为全双工特有收益。

DataTone 的歧义呈现、CLEO 的执行中介入和 VoCa 的语音画布协调均与本文直接相关。我们的增量需要落在分析对象来源与当前范围的具体分离方式，以及未来对照所产生的证据。仅换成 Olist 看板、增加时间戳或调用新模型，不足以构成这一增量。

\subsection{当前限制与复现}
系统目前收到新发话就撤销旧工作，尚不区分附和与分析修订；这限制了自然对话范围。原生模型历史未隔离，采样游标也未与词句自动对齐。正式对照前仍需冻结原生历史与可见支持差异，完成自然业务材料、独立意图评分及核听流程。界面可见不代表阅读，生成语音不代表理解。

本地验证使用脚本工具决策和测试音频；真实记录包含两个早期案例、八次诊断及另两次复测，没有参与者数据。各轮版本与时序不同，不能合并估计模型总体成功率或协作收益。预选历史证据也不能验证依赖播放内容的自然指代。复现材料同时保存版本、输入、图表状态、工具回执及原始失败；原始语音、数据和参与者资料按同意范围管理，不能随论文材料自动公开。

\section{结论}
\sys 面向全双工协同可视分析中的修订，将历史证据所指的对象与当前修改范围显式分开，并使用共享工具和可纠正的修订卡落实改动。系统区分提交、显示、播放和语义正确，避免把口头承诺当作分析完成。本地验证和真实诊断支持执行链的可实现性；旧八次四次最终状态正确，另两次复测一成一败。类型错误仍存在，旧竞态未被本轮真实触发，图像交付修复也仅有本地验证。所提出的双条件协议将进一步检验显式支持的正确性、理解收益与修复代价；在获得这些数据前，本文不宣称用户效果或全双工优越性。

\balance
\bibliographystyle{IEEEtran}
\bibliography{references}
\end{document}
